\section{Master consistency-strength map}
\label{sec:p06-consistency-strength}

The word \emph{stronger} is not used here without naming the comparison.
Theory inclusion, interpretability, conservation, ordinal analysis, and a
relative-consistency implication are different relations.  The table records
the exact relation in each row; the figure repeats only those labelled edges
and is not a single strict total order.

\begin{hbdefinition}{def:p06-pra-isigma-one}{The arithmetic theories \(\mathsf{PRA}\) and \(\mathsf I\Sigma_1\)}
Primitive recursive arithmetic \(\mathsf{PRA}\) is the first-order theory
with a symbol and defining equations for every primitive-recursive function
and induction for quantifier-free formulas.  The theory
\(\mathsf I\Sigma_1\) is \(\mathsf Q\) together with induction for every
\(\Sigma^0_1\) arithmetic formula, with parameters and universal closure as
in the PA convention.  These names are needed to identify the two lower
arithmetic calibrations in the table.  They stand as formal theories, differ
in language and induction presentation, and are used only under the
conservativity relation stated in the first row; their systematic treatment
comes in Part VII.
\end{hbdefinition}

\begingroup\footnotesize
\begin{longtable}{@{}p{0.18\textwidth}p{0.25\textwidth}p{0.28\textwidth}p{0.20\textwidth}@{}}
\caption{Master consistency- and proof-strength calibration.  Each row names
its comparison relation; no unlabelled strict consistency claim is intended.}
\label{tab:p06-consistency-strength-master}\\
\toprule
Track & Rungs & Supported comparison & Boundary and source \\
\midrule
\endfirsthead
\caption[]{Master consistency- and proof-strength calibration.  Each row names
its comparison relation; no unlabelled strict consistency claim is intended (continued)}\\
\toprule
Track & Rungs & Supported comparison & Boundary and source \\
\midrule
\endhead
\midrule\multicolumn{4}{r}{\emph{continued on next page}}\\
\endfoot
\endlastfoot
First-order arithmetic &
\(\mathsf{PRA}\), \(\mathsf I\Sigma_1\), \(\PA\) &
\(\mathsf I\Sigma_1\) is \(\Pi^0_2\)-conservative over the standard
presentation of \(\mathsf{PRA}\); they are equiconsistent and share the usual
\(\omega^\omega\) ordinal calibration.  The step from
\(\mathsf I\Sigma_1\) to \(\PA\) is strict in proof-theoretic strength. &
Parsons conservativity and ordinal-analysis calibration;
\cite{Parsons1970,Pohlers2009}. \\
\addlinespace
Reverse mathematics &
\(\mathsf{RCA}_0<\mathsf{WKL}_0<\mathsf{ACA}_0<
\mathsf{ATR}_0<\Pi^1_1\text{-}\mathsf{CA}_0<\mathsf Z_2\) &
Strict theorem inclusion in the common two-sorted language over the standard
base; the final rung is full second-order arithmetic. &
Definitions and representatives are in
Table~\ref{tab:p05-reverse-calibration}; \cite{Simpson2009}. \\
\addlinespace
Arithmetic bridge &
\(\PA\) and \(\mathsf{ACA}_0\) &
\(\mathsf{ACA}_0\) is conservative over \(\PA\) for first-order
arithmetical sentences; hence the displayed standard presentations are
equiconsistent. &
This is not literal theory inclusion because the languages differ;
\cite{Simpson2009}. \\
\addlinespace
Second-order-to-set bridge &
\(\Pi^1_1\text{-}\mathsf{CA}_0\), \(\mathsf Z_2\), \(\mathsf Z\) &
\(\mathsf Z\) proves that the standard two-sorted structure
\((\omega,\mathcal P(\omega))\) satisfies full second-order arithmetic;
therefore its consistency implies consistency of each displayed subsystem. &
An interpretation/standard-model implication, not a claimed adjacent strict
gap; \cite{Jech2003,Simpson2009}. \\
\addlinespace
Set-theory axioms &
\(\mathsf Z\), \(\mathsf{ZF}\), \(\mathsf{ZFC}\) &
\(\mathsf{ZF}\) extends \(\mathsf Z\) by Replacement.
\(\mathsf{ZFC}=\mathsf{ZF}+\mathsf{AC}\), while the constructible-inner-model
argument gives equiconsistency of \(\mathsf{ZF}\) and \(\mathsf{ZFC}\). &
No strict consistency separation between \(\mathsf{ZF}\) and
\(\mathsf{ZFC}\); \cite{Godel1938,Jech2003}. \\
\addlinespace
Large-cardinal extensions &
\(\mathsf{ZFC}+\) inaccessible, Mahlo, weakly compact, measurable,
supercompact &
Each stronger displayed cardinal property implies the lower properties in
Theorem~\ref{thm:p06-large-cardinal-implications}; an inaccessible
\(\kappa\) gives \((V_\kappa,\in)\entails\mathsf{ZFC}\). &
Conditional extensions only; no existence, adjacent strictness, or
exhaustiveness claim; \cite{Kanamori2003}. \\
\bottomrule
\end{longtable}
\endgroup

\begin{figure}[htbp]
\centering
\begin{tikzpicture}[
  x=1cm,y=.92cm,
  rung/.style={draw,rounded corners,align=center,font=\scriptsize,
               minimum width=2.75cm,minimum height=.52cm,inner sep=2pt},
  edge/.style={-{Stealth[length=2mm]},thick},
  bridge/.style={<->,>=Stealth,dashed,thick},
  note/.style={font=\tiny,align=center}
]
  \node[rung] (pra) at (-3.4,0) {\(\mathsf{PRA}\)};
  \node[rung] (is1) at (-3.4,1.1) {\(\mathsf I\Sigma_1\)};
  \node[rung] (pa) at (-3.4,2.2) {\(\PA\)};
  \draw[bridge] (pra) -- node[right,note]{\(\Pi^0_2\)-conservative\\
    equiconsistent} (is1);
  \draw[edge] (is1) -- node[right,note]{proof strength} (pa);

  \node[rung] (rca) at (3.4,0) {\(\mathsf{RCA}_0\)};
  \node[rung] (wkl) at (3.4,1.1) {\(\mathsf{WKL}_0\)};
  \node[rung] (aca) at (3.4,2.2) {\(\mathsf{ACA}_0\)};
  \node[rung] (atr) at (3.4,3.3) {\(\mathsf{ATR}_0\)};
  \node[rung] (pica) at (3.4,4.4)
    {\(\Pi^1_1\text{-}\mathsf{CA}_0\)};
  \node[rung] (z2) at (3.4,5.5) {\(\mathsf Z_2\)};
  \draw[edge] (rca) -- (wkl);
  \draw[edge] (wkl) -- (aca);
  \draw[edge] (aca) -- (atr);
  \draw[edge] (atr) -- (pica);
  \draw[edge] (pica) -- node[right,note]{theorem inclusion} (z2);
  \draw[bridge] (pa) -- node[above,note]{same first-order strength} (aca);

  \node[rung] (z) at (0,6.7) {\(\mathsf Z\)};
  \node[rung] (zfzfc) at (0,7.8)
    {\(\mathsf{ZF}\) / \(\mathsf{ZFC}\)};
  \node[rung] (inacc) at (0,8.9) {ZFC + inaccessible};
  \node[rung] (mahlo) at (0,10.0) {ZFC + Mahlo};
  \node[rung] (weak) at (0,11.1) {ZFC + weakly compact};
  \node[rung] (meas) at (0,12.2) {ZFC + measurable};
  \node[rung] (super) at (0,13.3) {ZFC + supercompact};
  \draw[edge,dashed] (z) -- node[right,note]{standard model of \(\mathsf Z_2\)} (z2);
  \draw[edge] (z) -- node[right,note]{add Replacement} (zfzfc);
  \draw[edge] (inacc) -- node[right,note]{\(V_\kappa\models\mathsf{ZFC}\)} (zfzfc);
  \draw[edge] (mahlo) -- node[right,note]{cardinal implication} (inacc);
  \draw[edge] (weak) -- node[right,note]{cardinal implication} (mahlo);
  \draw[edge] (meas) -- node[right,note]{cardinal implication} (weak);
  \draw[edge] (super) -- node[right,note]{cardinal implication} (meas);
  \node[note,text width=4.2cm] at (-3.2,7.1)
    {ZF and ZFC are equiconsistent via the constructible inner model; the
     shared box is not a claim that Choice is provable in ZF.};
\end{tikzpicture}
\caption{Consistency- and proof-strength map. Upward inclusion and
proof-strength arrows, downward cardinal-property and \(V_\kappa\)-model
arrows, and dashed standard-model or equiconsistency bridges record the
different relations named in Table~\ref{tab:p06-consistency-strength-master}.
The diagram
does not assert one strict total order, large-cardinal existence, strictness of
adjacent large-cardinal consistency strengths, or exhaustion of stronger
theories.}
\label{fig:p06-consistency-strength}
\end{figure}
